\section{An open compliance benchmark}\label{sec:benchmark}

The crosswalk turns five regulatory texts into one structured object. The benchmark turns that object into a \emph{measurement}: given a machine-learning pipeline (a repository, its artefacts, and its documentation), how much of the union standard does it satisfy, and which regime does it fall short of?

\subsection{Scoring model}\label{sec:scoring}

Each of the twelve requirements of \Cref{tab:crosswalk} is scored on a pipeline as \emph{met}, \emph{partial}, or \emph{absent}. A requirement is \emph{met} only when the pipeline carries the artefact the strictest regime in that row demands, so for record retention (row 7), a pipeline is scored against the EU AI Act's $\geq$6-month floor even when deployed where no floor exists. From the twelve per-requirement scores we report two aggregates:

\begin{itemize}[leftmargin=*]
    \item a \textbf{per-regime score}: the fraction of that regime's applicable requirements the pipeline meets, so a firm can read its standing under SR 11-7, the EU AI Act, PRA SS1/23, MAS FEAT, or OSFI E-23 individually; and
    \item a \textbf{union score}: the fraction of the full twelve met at the strictest level, the single number that certifies a pipeline as deployable across all five jurisdictions without per-market rework.
\end{itemize}

Six of the twelve requirements are mechanically checkable from a repository using the same static-analysis approach validated in the predecessor rubric of \citet{KhanMrAudit2026}: inventory metadata (row 1), documentation and environment (row 2), data references (row 3), version control (row 4), the presence of a retention policy (row 7), and a decommissioning record (row 12). The remaining six require either code execution or manual review against a published codebook; the scoring harness automates the first set and structures the second.

\subsection{What is released}\label{sec:released}

The artefact accompanying this paper is released under permissive licences (CC-BY-4.0 for the crosswalk, MIT for code) at the companion repository and comprises: (i) the crosswalk of \Cref{tab:crosswalk} as a machine-readable table (\texttt{crosswalk.csv}), keyed by requirement and regime with the governing provision in each cell; (ii) a reference scoring rubric that operationalises each requirement as a checkable predicate; and (iii) a reference scoring harness with a public leaderboard, seeded on public-data machine-learning-finance pipelines so that the first entries are reproducible by any reader.

We are explicit about the version state: this is a v0.1 seed. The crosswalk is complete for the seven regimes studied; the harness automates the six mechanically checkable requirements and ships the codebook for the manual six; and the seed leaderboard is intentionally small and community-extensible rather than a comprehensive ranking. The benchmark is designed to be \emph{extended} (by adding regimes as columns, pipelines as rows, and updated provisions as the law changes), not to be a finished league table.

\subsection{Why a benchmark, not just a survey}\label{sec:why_benchmark}

A static survey of compliance goes out of date the moment a regulation changes. A versioned, open benchmark stays current: when OSFI E-23 takes effect in 2027, or when the U.S.\ revised guidance settles, a column is updated and every pipeline is re-scored without re-reading five texts. The benchmark also makes compliance \emph{comparable} across pipelines, which a prose survey cannot, and gives practitioners something to engineer against instead of a document to interpret.
